%%%%%%%%%%%%%%%%%%%%%%%%%%%%%%%%%%%%%%%%%%%%%%%%%%%%%%%%%%%
% --------------------------------------------------------
%  _____           
% |_   _|_ _ _   _ 
%   | |/ _` | | | |
%   | | (_| | |_| |
%   |_|\__,_|\__,_|
%
% LaTeX2e Template
% Version 2.6.0 (21/04/2026)
%
% Author: 
% Guillermo Jimenez (memo.notess1@gmail.com)
% 
% License:
% Creative Commons CC BY 4.0
% --------------------------------------------------------
%%%%%%%%%%%%%%%%%%%%%%%%%%%%%%%%%%%%%%%%%%%%%%%%%%%%%%%%%%%
% --------------------------------------------------------
% Github Repository:
% https://github.com/MemoJimenez/Tau-class
% --------------------------------------------------------
%%%%%%%%%%%%%%%%%%%%%%%%%%%%%%%%%%%%%%%%%%%%%%%%%%%%%%%%%%%

\documentclass[9pt,a4paper,twocolumn,twoside]{tau-class/tau}
\usepackage[utf8]{inputenc}
\usepackage[english]{babel}
\usepackage{kotex}
\ifdefined\setmainhangulfont
\setmainhangulfont{Noto Serif CJK KR}[BoldFont = Noto Serif CJK KR Bold]
\fi
\usepackage{biblatex}
\usepackage{hyperref}   % makes references clickable
\usepackage{cleveref}   % load AFTER hyperref
\usepackage{amsmath}
\usepackage{tikz}
\usepackage{float}


\newcommand{\dotlink}{\mathrel{\bullet\joinrel\relbar\joinrel\bullet}}

% Spanish babel recommendation
% \usepackage[spanish,es-nodecimaldot,es-noindentfirst]{babel} 

%----------------------------------------------------------
% Title & Authors/Affiliations
%----------------------------------------------------------

\doctype{Manual}
\title{Basics of Hangul the Korean Writing System: \\ A Guide for Beginners}

%----------------------------------------------------------

\author[a,1]{Hyungi Choi\thanks{\href{mailto:choihyungi92@gmail.com}{choihyungi92@gmail.com}}}


%----------------------------------------------------------
% Document-, Footer- information
%----------------------------------------------------------

\dates{This manuscript was compiled on \today}

\docinfo{} 

%----------------------------------------------------------

\footinfo{personal Miniproject}
\theday{\today}
\organization{Vienna University of Technology, Austria}
\leadauthor{Hyungi Choi}

%----------------------------------------------------------
% Abstract and Keywords
%----------------------------------------------------------

\begin{abstract}    
    This document serves as a comprehensive guide for beginners to learn the basics of Hangul, the Korean writing system. It covers the fundamental principles of Hangul, mainly its structure, and pronunciation. The guide is designed to provide learners with a solid foundation in reading and writing Hangul.
\end{abstract}

%----------------------------------------------------------

%\keywords{tau-class, research article, lab report, academic writing}

%----------------------------------------------------------
\begin{document}

    \maketitle
    \thispagestyle{firststyle}
    \tableofcontents

\section{Introduction}
    Contrary to a common misconception, Hangul is not a logographic writing system like Chinese characters. Instead, it is a featural alphabet consisting of 14 basic consonants and 10 basic vowels written in blocks to represent the syllables of the Korean language. It was created in 1443 by King Sejong the Great and his scholars with the purpose of promoting literacy among the common people, therefore it is characterized by its simplicity and ease of learning~\cite{hunminjeongeum}. 

\section{Principles}
    It is crucial to understand the basic principles of Hangul before diving into its structure and pronunciation. The Korean writing system is based on a specific arrangement of consonants and vowels, that result in 2-dimensional blocks representing syllables, which are then written linearly to form words. This unusual arrangement of letters is probably the most challenging aspect of learning Hangul.
    
    The consonants are designed to represent the shape of the speech organs used to produce them, while the vowels are designed to represent the philosophical concepts of the universe at the time of its creation~\cite{hunminjeongeum}. While the visual appeal of combining letters into single syllabic blocks may vary from person to person, it cannot be denied that the clear visual distinction between consonants and vowels is a welcome sign for any reader.

\section{Structure}
    Every syllable is written as a single block built from up to three parts: an initial consonant $\mathrm{c}_i$, a total vowel $\mathrm{v}_t$, and a final optional consonant $\mathrm{c}_o$.
    
    \begin{figure}[H]
        \centering
        \begin{tikzpicture}[scale=0.9]
            \draw (0,2) rectangle (2,3) node[pos=.5] {$c_1$};
            \draw (0,1) rectangle (2,2) node[pos=.5] {$v_1$};
            \draw (2,1) rectangle (3,3) node[pos=.5] {$v_2$};
            \draw (0,0) rectangle (1.5,1) node[pos=.5] {$c_2$};
            \draw (1.5,0) rectangle (3,1) node[pos=.5] {$c_3$};
            \node[anchor=west, align=left] at (3.4,1.5) {$c$\ldots consonant\\ $v$\ldots vowel};
        \end{tikzpicture}
        \caption{Basic structure of a Hangul syllable block}
        \label{fig:block}
    \end{figure}

    \cref{fig:block} shows the block structure. Letters are read in this order:

    \begin{equation}
        \mathrm{Syllable} = 
        \underbrace{\mathrm{c}_1}_{\text{initial consonant}\ \mathrm{c}_i} + \underbrace{(\mathrm{v}_1 + \mathrm{v}_2)}_{\text{total vowel}\ \mathrm{v}_t} + \underbrace{(\mathrm{c}_2 + \mathrm{c}_3)}_{\text{optional consonant}\ \mathrm{c}_o}  \label{eq:syllable_structure}
    \end{equation}
    $\mathrm{c}_o$ \cref{eq:syllable_structure} If the optional consonant $\mathrm{c}_o$ uses both $\mathrm{c}_2$ and $\mathrm{c}_3$, then they are spoken together as a single consonant sound by compressing the sound of the individual consonants.
\section{Consonants}
    The consonants are grouped by the speech organ that produces them. Each consonant is shown in this notation:

    \begin{equation}
        \begin{pmatrix}
            \multicolumn{3}{c}{\textbf{\Large hangul letter}} \\
            \text{$\text{sound}^{\text{default}}_{\text{reference word}}$} & \text{$\text{sound}^{\text{special case Nr.}}_{\text{reference word}}(\text{trigger})$} & \text{as $\textbf{\large c}_\mathbf{1}$}\\
            \text{$\text{sound}^{\text{default}}_{\text{reference word}}$} & \text{$\text{sound}^{\text{special case Nr.}}_{\text{reference word}}(\text{trigger})$} & \text{as $\textbf{\large c}_\mathbf{2}$}\\
            \text{$\text{sound}^{\text{default}}_{\text{reference word}}$} & \text{$\text{sound}^{\text{special case Nr.}}_{\text{reference word}}(\text{trigger})$} & \text{as $\textbf{\large c}_\mathbf{3}$}\\
        \end{pmatrix}
    \end{equation}
   
    Pronunciations are simplified for beginners, based on the official rules~\cite{pronunciation_standard}, and written as an English speaker would read them rather than in the official romanization (e.g. ㅡ as “eww” instead of “eu”). Onomatopoeia and interjections are marked with quotes, $\varnothing$ means silent, and numbered special cases are listed below each bracket.

    \subsection{Tongue, velar}
        Back of tongue against the soft palate; ㄱ shows the tongue root.

        \begin{equation}
            \begin{pmatrix}
                \pmb{\text{\Large ㄱ}}\\ 
                \mathrm{g}_\mathrm{\textbf{g}o}\\ 
                \mathrm{g}_\mathrm{\textbf{g}o}\\
                \mathrm{g}_\mathrm{\textbf{g}o}
            \end{pmatrix}
            \dotlink
            \begin{pmatrix}
                \pmb{\text{\Large ㅋ}} \\ 
                \mathrm{k}_\mathrm{\textbf{k}ite}\\ 
                \mathrm{g}_\mathrm{\textbf{g}o}\\
                \mathrm{-}
            \end{pmatrix}
            \dotlink
             \begin{pmatrix}
                \pmb{\text{\Large ㄲ}} \\
                \mathrm{gg}_\mathrm{''\textbf{G}ulp!''}\\
                \mathrm{-}\\
                \mathrm{-}
            \end{pmatrix}
            \label{eq:giyeok}
        \end{equation}
       
    
    \subsection{Tongue, alveolar}
        Tongue tip touches the ridge behind the upper teeth; ㄴ shows the tongue.

        \begin{equation}
            \begin{pmatrix}
                \pmb{\text{\Large ㄴ}}\\
                \mathrm{n}_\mathrm{\textbf{n}o}\\
                \mathrm{n}_\mathrm{\textbf{n}o}\\
                \mathrm{-}
            \end{pmatrix}
            \dotlink
            \begin{pmatrix}
                \pmb{\text{\Large ㄷ}}\\
                \mathrm{d}_\mathrm{\textbf{d}o}\\
                \mathrm{d}_\mathrm{\textbf{d}o}\\
                \mathrm{-}
            \end{pmatrix}
            \dotlink
            \begin{pmatrix}
                \pmb{\text{\Large ㅌ}}\\
                \mathrm{t}_\mathrm{\textbf{t}op}\\
                \mathrm{d}_\mathrm{\textbf{d}o}\\
                \mathrm{t}_\mathrm{\textbf{t}op}
            \end{pmatrix}
            \dotlink
            \begin{pmatrix}
                \pmb{\text{\Large ㄸ}}\\
                \mathrm{dd}_\mathrm{''\textbf{D}ing!''}\\
                \mathrm{-}\\
                \mathrm{-}
            \end{pmatrix}
            \label{eq:nieun}
        \end{equation}
        \begin{equation}
            \begin{pmatrix}
                \multicolumn{2}{c}{\pmb{\text{\Large ㄹ}}}\\
                \mathrm{r}_\mathrm{Ma\textbf{r}ia} & \mathrm{l}^{~\ref{item:rieut_special_case_1}}_\mathrm{ba\textbf{ll}}(\mathrm{c}_{end}^{\text{prev}} = \text{ㄹ})\\
                \mathrm{l}_\mathrm{ba\textbf{ll}}\\
                \mathrm{-}
            \end{pmatrix}
            \label{eq:rieut}
        \end{equation}
        
        ㄹ \cref{eq:rieut} special cases: 
        \begin{enumerate}
            \item \textbf{If previous syllable ends in} \pmb{ㄹ}, then \pmb{ㄹ} in $\textbf{c}_1$ is read as \textbf{l not r}.\\ (빨\textbf{리} $\rightarrow$ bbal-\textbf{r}i $\pmb\times$, bbal-\textbf{l}i $\checkmark$ ). \label{item:rieut_special_case_1}
        \end{enumerate}
    
    \subsection{Lips, bilabial}
        Lips pressed together; ㅁ shows the mouth.

        \begin{equation}
            \begin{pmatrix}
                \pmb{\text{\Large ㅁ}}\\
                \mathrm{m}_\mathrm{\textbf{m}oon}\\
                \mathrm{m}_\mathrm{\textbf{m}oon}\\
                \mathrm{m}_\mathrm{\textbf{m}oon}
            \end{pmatrix}
            \dotlink
            \begin{pmatrix}
                \pmb{\text{\Large ㅂ}}\\
                \mathrm{b}_\mathrm{\textbf{b}oy}\\
                \mathrm{b}_\mathrm{\textbf{b}oy}\\
                \mathrm{b}_\mathrm{\textbf{b}oy}\\
            \end{pmatrix}
            \dotlink
            \begin{pmatrix}
                \pmb{\text{\Large ㅍ}}\\
                \mathrm{p}_\mathrm{\textbf{p}ie}\\
                \mathrm{b}_\mathrm{\textbf{b}oy}\\
                \mathrm{b}_\mathrm{\textbf{b}oy}
            \end{pmatrix}
            \dotlink
            \begin{pmatrix}
                \pmb{\text{\Large ㅃ}}\\
                \mathrm{bb}_\mathrm{''\textbf{B}ang!''}\\
                -\\
                -
            \end{pmatrix}
            \label{eq:mieum}
        \end{equation}
    
    \subsection{Teeth, sibilant}
        Air hisses past the teeth; ㅅ shows the gap between teeth.

       \begin{equation}
            \begin{pmatrix}
                \multicolumn{2}{c}{\pmb{\text{\Large ㅅ}}}\\
                \mathrm{s}_\mathrm{\textbf{s}un} & \mathrm{sh}_\mathrm{\textbf{sh}e}^{~\ref{item: shieut_special_case_1}}(\mathrm{v}_t = \text{ㅣ})\\
                \mathrm{d}_\mathrm{\textbf{d}o} & \mathrm{s}_\mathrm{\textbf{s}un}^{~\ref{item: shieut_special_case_2}}(\mathrm{c}^{\text{next}}_1 = \text{ㅇ})\\
                \mathrm{d}_\mathrm{\textbf{d}o} & \mathrm{s}_\mathrm{\textbf{s}un}^{~\ref{item: shieut_special_case_2}}(\mathrm{c}^{\text{next}}_1 = \text{ㅇ})\\
                           
            \end{pmatrix}
            \dotlink
            \begin{pmatrix}
                \pmb{\text{\Large ㅆ}}\\
                \mathrm{ss}_\mathrm{\textbf{s}ea}\\
                -\\
                -
            \end{pmatrix}
            \label{eq:shieut}
        \end{equation}

        ㅅ \cref{eq:shieut} special cases:
        \begin{enumerate}
            \item \textbf{If $\mathrm{v}_1$ is empty and $\mathrm{v}_2$ is} \pmb{ㅣ}, then \pmb{ㅅ} in $\textbf{c}_1$is read as \textbf{sh not s}\\ 
            (시 → \textbf{s}i \pmb{$\times$}, \textbf{sh}i $\checkmark$). 
            \label{item: shieut_special_case_1}
            \item \textbf{If the next syllable starts with} \pmb{ㅇ}, then any \pmb{ㅅ} in $\textbf{c}_2$ and $\textbf{c}_3$ is read as \textbf{s not d}.\\ 
            (\textbf{옷}이 → o\textbf{d}-i \pmb{$\times$}, o\textbf{s}-i $\checkmark$)\\ 
            (\textbf{있}어 → id\textbf{d}-uh \pmb{$\times$}, is\textbf{s}-uh $\checkmark$). 
            \label{item: shieut_special_case_2}
        \end{enumerate}

        \begin{equation}
            \begin{pmatrix}
                \multicolumn{2}{c}{\pmb{\text{\Large ㅈ}}}\\
                \mathrm{j}_\mathrm{\textbf{j}am}\\
                \mathrm{d}_\mathrm{\textbf{d}o} &  \mathrm{j}_\mathrm{\textbf{j}am}^{~\ref{item: jieut_special_case_1}}(\mathrm{c}^{\text{next}}_1 = \text{ㅇ})\\
                \mathrm{d}_\mathrm{\textbf{d}o} &  \mathrm{j}_\mathrm{\textbf{j}am}^{~\ref{item: jieut_special_case_1}}(\mathrm{c}^{\text{next}}_1 = \text{ㅇ})\\
            \end{pmatrix}
            \dotlink
            \begin{pmatrix}
                \pmb{\text{\Large ㅊ}}\\
                \mathrm{ch}_\mathrm{\textbf{ch}ild}\\
                \mathrm{d}_\mathrm{\textbf{d}o}\\
                -
            \end{pmatrix}
            \dotlink
            \begin{pmatrix}
                \pmb{\text{\Large ㅉ}}\\
                \mathrm{jj}_\mathrm{Gu\textbf{cc}i}\\
                -\\
                -
            \end{pmatrix}
            \label{eq:jieut}
        \end{equation}
        
        ㅈ \cref{eq:jieut} special cases:
        \begin{enumerate}
            \item \textbf{If the next syllable starts with} \pmb{ㅇ}, then any \pmb{ㅈ} in $\textbf{c}_2$ and $\textbf{c}_3$ is read as \textbf{j not d}.\\ 
            (\textbf{맞}아 → ma\textbf{d}-a \pmb{$\times$}, ma\textbf{j}-a $\checkmark$)\\ 
            (\textbf{앉}아 → an\textbf{d}-a \pmb{$\times$}, an\textbf{j}-a $\checkmark$)
            \label{item: jieut_special_case_1}
        \end{enumerate}

       
    
    \subsection{Throat, glottal}
       Sound from the throat; ㅇ shows the open throat

       \begin{equation}
            \begin{pmatrix}
                \pmb{\text{\Large ㅇ}}\\
                \mathrm{\varnothing}\\
                \mathrm{ng}_\mathrm{si\textbf{ng}}\\
                -
            \end{pmatrix}
            \dotlink
            \begin{pmatrix}
                \pmb{\text{\Large ㅎ}}\\
                \mathrm{h}_\mathrm{\textbf{h}at}\\
                \mathrm{\varnothing}\\
                \varnothing
            \end{pmatrix}
            \label{eq:ieung}
       \end{equation}
 
\section{Vowels}
    The vowels are built from $\cdot$, ㅡ, and ㅣ (see the introduction). For reading, it is easiest to sort them by where they sit in the block: horizontal ($\mathrm{v}_1$), vertical ($\mathrm{v}_2$), and compound ($\mathrm{v}_1 + \mathrm{v}_2$).
    Each vowel is shown in this notation:

    \begin{equation}
        \begin{pmatrix}
            \text{hangul letter}\\
            \mathrm{vowel}_\text{reference word}
        \end{pmatrix}
    \end{equation}

    \subsection{Horizontal Vowels $\mathrm{v}_1$}
        Based on ㅡ, written below $\mathrm{c}_1$.
         \begin{equation}
            \begin{pmatrix}
                \pmb{\text{\Large ㅡ}}\\
                \mathrm{eww}_\mathrm{''\textbf{eww!}''}
            \end{pmatrix}
            \label{eq:eu}
        \end{equation}

        \begin{equation}
            \begin{pmatrix}
                \pmb{\text{\Large ㅗ}}\\
                \mathrm{o}_\mathrm{g\textbf{o}}
            \end{pmatrix}
            \dotlink
            \begin{pmatrix}
                \pmb{\text{\Large ㅛ}}\\
                \mathrm{yo}_\mathrm{\textbf{yo}ga}
            \end{pmatrix}
            \label{eq:o}
        \end{equation}

        \begin{equation}
            \begin{pmatrix}
                \pmb{\text{\Large ㅜ}}\\
                \mathrm{oo}_\mathrm{m\textbf{oo}n}
            \end{pmatrix}
            \dotlink
            \begin{pmatrix}
                \pmb{\text{\Large ㅠ}}\\
                \mathrm{you}_\mathrm{\textbf{you}}
            \end{pmatrix}
            \label{eq:u}
        \end{equation}
       
    \subsection{Vertical Vowels $\mathrm{v}_2$}
        Based on ㅣ, written to the right of $\mathrm{c}_1$.
        
        \begin{equation}
            \begin{pmatrix}
                \pmb{\text{\Large ㅣ}}\\
                \mathrm{i}_\mathrm{s\textbf{ee}}
            \end{pmatrix}
            \label{eq:i}
        \end{equation}

        \begin{equation}
            \begin{pmatrix}
                \pmb{\text{\Large ㅏ}}\\
                \mathrm{a}_\mathrm{sp\textbf{a}}
            \end{pmatrix}
            \dotlink
            \begin{pmatrix}
                \pmb{\text{\Large ㅑ}}\\
                \mathrm{ya}_\mathrm{\textbf{ya}rd}
            \end{pmatrix}
            \label{eq:a}
        \end{equation}

        \begin{equation}
            \begin{pmatrix}
                \pmb{\text{\Large ㅓ}}\\
                \mathrm{uh}_\mathrm{s\textbf{u}n}
            \end{pmatrix}
            \dotlink
            \begin{pmatrix}
                \pmb{\text{\Large ㅕ}}\\
                \mathrm{yuh}_\mathrm{\textbf{you}ng}
            \end{pmatrix}
            \label{eq:eo}
        \end{equation}

        \begin{equation}
            \begin{pmatrix}
                \pmb{\text{\Large ㅐ}}\\
                \mathrm{ae}_\mathrm{c\textbf{a}t}
            \end{pmatrix}
            \dotlink
            \begin{pmatrix}
                \pmb{\text{\Large ㅒ}}\\
                \mathrm{yae}_\mathrm{\textbf{yea}h}
            \end{pmatrix}
            \label{eq:ae}
        \end{equation}

        \begin{equation}
            \begin{pmatrix}
                \pmb{\text{\Large ㅔ}}\\
                \mathrm{e}_\mathrm{h\textbf{e}y}
            \end{pmatrix}
            \dotlink
            \begin{pmatrix}
                \pmb{\text{\Large ㅖ}}\\
                \mathrm{ye}_\mathrm{\textbf{ya}y}
            \end{pmatrix}
            \label{eq:e}
        \end{equation}

\subsection{Compound vowels $\mathrm{v}_1 + \mathrm{v}_2$}
    A horizontal and a vertical vowel combined, filling both positions.
    
    \begin{equation}
       \begin{pmatrix}
            \pmb{\text{\Large ㅢ}}\\
            \mathrm{ui}_\mathrm{''\textbf{eww!}'' + \textbf{ee}}
        \end{pmatrix}
        \label{eq:ui}
    \end{equation}

    \begin{equation}
        \begin{pmatrix}
            \pmb{\text{\Large ㅚ}}\\
            \mathrm{oe}_\mathrm{\textbf{we}t}
        \end{pmatrix}
        \dotlink
        \begin{pmatrix}
            \pmb{\text{\Large ㅘ}}\\
            \mathrm{wa}_\mathrm{\textbf{wa}sabi}
        \end{pmatrix}
        \dotlink
        \begin{pmatrix}
            \pmb{\text{\Large ㅙ}}\\
            \mathrm{wae}_\mathrm{\textbf{wa}gon}
        \end{pmatrix}
        \label{eq:oe}
    \end{equation}

    \begin{equation}
        \begin{pmatrix}
            \pmb{\text{\Large ㅟ}}\\
            \mathrm{wi}_\mathrm{\textbf{wee}k}
        \end{pmatrix}
        \dotlink
        \begin{pmatrix}
            \pmb{\text{\Large ㅝ}}\\
            \mathrm{wo}_\mathrm{\textbf{wo}nder}
        \end{pmatrix}
        \dotlink
        \begin{pmatrix}
            \pmb{\text{\Large ㅞ}}\\
            \mathrm{we}_\mathrm{\textbf{we}dding}
        \end{pmatrix}
        \label{eq:wi}
    \end{equation}

\section{Examples}
    Each example breaks a syllable into its parts and then reads them in order, using this notation:

    \begin{equation}
        \begin{pmatrix} \pmb{\text{\Large hangul letter}} \\ \boxed{\text{position in Block}} \\ \text{default sound} \\ \text{special case}\end{pmatrix}
    \end{equation}

    \subsection{Initial consonant + single vowel}
    \begin{align}
        \textbf{\Huge 가} &= 
        \underbrace{
            \begin{pmatrix}\pmb{\text{\Large ㄱ}}\\\boxed{\mathrm{c}_1}\\ \mathrm{g}\\ \mathrm{} \end{pmatrix}
        }_{\mathrm{c}_i}
        +
        \underbrace{
            \begin{pmatrix}             
                \pmb{\text{\Large }}&\pmb{\text{\Large ㅏ}}\\
                \boxed{\mathrm{v}_1}&\boxed{\mathrm{v}_2}\\
                \mathrm{}&\mathrm{a}\\
                \mathrm{}&\mathrm{}
            \end{pmatrix}
        }_{\mathrm{v}_t}
        +
        \underbrace{
            \begin{pmatrix}
                \pmb{\text{\Large }} & \pmb{\text{\Large }}\\
                \boxed{\mathrm{c}_2} & \boxed{\mathrm{c}_3}\\
                \mathrm{} & \mathrm{}\\
                \mathrm{} & \mathrm{}
            \end{pmatrix}
        }_{\mathrm{c}_o}
        \\&=
        \begin{pmatrix}
            \pmb{\text{\Large ㄱ}} & \pmb{\text{\Large ㅏ}} & \pmb{\text{\Large }}\\
            \boxed{\mathrm{c}_i} & \boxed{\mathrm{v}_t} & \boxed{\mathrm{c}_o}\\
            \mathrm{g} & \mathrm{a} & \mathrm{}\\
            \mathrm{} & \mathrm{} & \mathrm{}
        \end{pmatrix}
        =
        \begin{pmatrix}\pmb{\text{\Large 가}}\\\boxed{\mathrm{syllable}}\\\mathrm{ga}\\\mathrm{}\end{pmatrix}
    \end{align}\\
    
    \begin{align}
        \textbf{\Huge 고} &= 
        \underbrace{
            \begin{pmatrix}\pmb{\text{\Large ㄱ}}\\\boxed{\mathrm{c}_1}\\\mathrm{g}\\\mathrm{}\end{pmatrix}
        }_{\mathrm{c}_i}
        +
        \underbrace{
            \begin{pmatrix}             
                \pmb{\text{\Large ㅗ}} & \pmb{\text{\Large }}\\
                \boxed{\mathrm{v}_1} & \boxed{\mathrm{v}_2}\\
                \mathrm{o} & \mathrm{}\\
                \mathrm{} & \mathrm{}
                \end{pmatrix}
        }_{\mathrm{v}_t}
            +
        \underbrace{
            \begin{pmatrix}
                \pmb{\text{\Large }} & \pmb{\text{\Large }}\\
                \boxed{\mathrm{c}_2} & \boxed{\mathrm{c}_3}\\
                \mathrm{} & \mathrm{}\\
                \mathrm{} & \mathrm{}
            \end{pmatrix}
        }_{\mathrm{c}_o}
        \\&=
        \begin{pmatrix}
            \pmb{\text{\Large ㄱ}} & \pmb{\text{\Large ㅗ}} & \pmb{\text{\Large }}\\
            \boxed{\mathrm{c}_i} & \boxed{\mathrm{v}_t} & \boxed{\mathrm{c}_o}\\
            \mathrm{g} & \mathrm{o} & \mathrm{}\\
            \mathrm{} & \mathrm{} & \mathrm{}
        \end{pmatrix}
        =
        \begin{pmatrix}\pmb{\text{\Large 고}}\\\boxed{\mathrm{syllable}}\\\mathrm{go}\\\mathrm{}\end{pmatrix}
    \end{align}
        \subsection{Initial consonant + compound vowel}
        \begin{align}
        \textbf{\Huge 와} &= 
        \underbrace{
            \begin{pmatrix}\pmb{\text{\Large ㅇ}}\\\boxed{\mathrm{c}_1}\\\mathrm{\varnothing}\\\mathrm{}\end{pmatrix}
        }_{\mathrm{c}_i}
        +
        \underbrace{
            \begin{pmatrix}             
                \pmb{\text{\Large ㅗ}} & \pmb{\text{\Large ㅏ}}\\
                \boxed{\mathrm{v}_1} & \boxed{\mathrm{v}_2}\\
                \mathrm{o} & \mathrm{a}\\
                \mathrm{} & \mathrm{}
                \end{pmatrix}
        }_{\mathrm{v}_t}
            +
        \underbrace{
            \begin{pmatrix}
                \pmb{\text{\Large }} & \pmb{\text{\Large }}\\
                \boxed{\mathrm{c}_2} & \boxed{\mathrm{c}_3}\\
                \mathrm{} & \mathrm{}\\
                \mathrm{} & \mathrm{}
            \end{pmatrix}
        }_{\mathrm{c}_o}
        \\&=
        \begin{pmatrix}
            \pmb{\text{\Large ㅇ}} & \pmb{\text{\Large ㅘ}} & \pmb{\text{\Large }}\\
            \boxed{\mathrm{c}_i} & \boxed{\mathrm{v}_t} & \boxed{\mathrm{c}_o}\\
            \mathrm{\varnothing} & \mathrm{wa} & \mathrm{}\\
            \mathrm{} & \mathrm{} & \mathrm{}
        \end{pmatrix}
        =
        \begin{pmatrix}\pmb{\text{\Large 와}}\\\boxed{\mathrm{syllable}}\\\mathrm{wa}\\\mathrm{}\end{pmatrix}
    \end{align}
    \subsection{Initial consonant + single vowel + optional consonant}
    \begin{align}
        \textbf{\Huge 밥} &= 
        \underbrace{
            \begin{pmatrix}\pmb{\text{\Large ㅂ}}\\\boxed{\mathrm{c}_1}\\\mathrm{b}\\\mathrm{}\end{pmatrix}
        }_{\mathrm{c}_i}
        +
        \underbrace{
            \begin{pmatrix}             
                \pmb{\text{\Large }} & \pmb{\text{\Large ㅏ}}\\
                \boxed{\mathrm{v}_1} & \boxed{\mathrm{v}_2}\\
                \mathrm{} & \mathrm{a}\\
                \mathrm{} & \mathrm{}
                \end{pmatrix}
        }_{\mathrm{v}_t}
            +
        \underbrace{
            \begin{pmatrix}
                \pmb{\text{\Large ㅂ}} & \pmb{\text{\Large }}\\
                \boxed{\mathrm{c}_2} & \boxed{\mathrm{c}_3}\\
                \mathrm{b} & \mathrm{}\\
                \mathrm{} & \mathrm{}
            \end{pmatrix}
        }_{\mathrm{c}_o}
        \\&=
        \begin{pmatrix}
            \pmb{\text{\Large ㅂ}} & \pmb{\text{\Large ㅏ}} & \pmb{\text{\Large ㅂ}}\\
            \boxed{\mathrm{c}_i} & \boxed{\mathrm{v}_t} & \boxed{\mathrm{c}_o}\\
            \mathrm{b} & \mathrm{a} & \mathrm{b}\\
            \mathrm{} & \mathrm{} & \mathrm{}
        \end{pmatrix}
        =
        \begin{pmatrix}\pmb{\text{\Large 밥}}\\\boxed{\mathrm{syllable}}\\\mathrm{bab}\\\mathrm{}\end{pmatrix}
    \end{align}\\
    
    \begin{align}
        \textbf{\Huge 꽃} &= 
        \underbrace{
            \begin{pmatrix}\pmb{\text{\Large ㄲ}}\\\boxed{\mathrm{c}_1}\\\mathrm{gg}\\\mathrm{}\end{pmatrix}
        }_{\mathrm{c}_i}
        +
        \underbrace{
            \begin{pmatrix}             
                \pmb{\text{\Large ㅗ}} & \pmb{\text{\Large }}\\
                \boxed{\mathrm{v}_1} & \boxed{\mathrm{v}_2}\\
                \mathrm{o} & \mathrm{}\\
                \mathrm{} & \mathrm{}
                \end{pmatrix}
        }_{\mathrm{v}_t}
            +
        \underbrace{
            \begin{pmatrix}
                \pmb{\text{\Large ㅊ}} & \pmb{\text{\Large }}\\
                \boxed{\mathrm{c}_2} & \boxed{\mathrm{c}_3}\\
                \mathrm{d} & \mathrm{}\\
                \mathrm{} & \mathrm{}
            \end{pmatrix}
        }_{\mathrm{c}_o}
        \\&=
        \begin{pmatrix}
            \pmb{\text{\Large ㄲ}} & \pmb{\text{\Large ㅗ}} & \pmb{\text{\Large ㅊ}}\\
            \boxed{\mathrm{c}_i} & \boxed{\mathrm{v}_t} & \boxed{\mathrm{c}_o}\\
            \mathrm{gg} & \mathrm{o} & \mathrm{d}\\
            \mathrm{} & \mathrm{} & \mathrm{}
        \end{pmatrix}
        =
        \begin{pmatrix}\pmb{\text{\Large 꽃}}\\\boxed{\mathrm{syllable}}\\\mathrm{ggod}\\\mathrm{}\end{pmatrix}
    \end{align}
    \subsection{Initial consonant + compound vowel + optional consonant}
    \begin{align}
        \textbf{\Huge 왔} &= 
        \underbrace{
            \begin{pmatrix}\pmb{\text{\Large ㅇ}}\\\boxed{\mathrm{c}_1}\\\mathrm{\varnothing}\\\mathrm{}\end{pmatrix}
        }_{\mathrm{c}_i}
        +
        \underbrace{
            \begin{pmatrix}             
                \pmb{\text{\Large ㅗ}} & \pmb{\text{\Large ㅏ}}\\
                \boxed{\mathrm{v}_1} & \boxed{\mathrm{v}_2}\\
                \mathrm{o} & \mathrm{a}\\
                \mathrm{} & \mathrm{}
                \end{pmatrix}
        }_{\mathrm{v}_t}
            +
        \underbrace{
            \begin{pmatrix}
                \pmb{\text{\Large ㅅ}} & \pmb{\text{\Large ㅅ}}\\
                \boxed{\mathrm{c}_2} & \boxed{\mathrm{c}_3}\\
                \mathrm{d} & \mathrm{d}\\
                \mathrm{} & \mathrm{}
            \end{pmatrix}
        }_{\mathrm{c}_o}
        \\&=
        \begin{pmatrix}
            \pmb{\text{\Large ㅇ}} & \pmb{\text{\Large ㅘ}} & \pmb{\text{\Large ㅆ}}\\
            \boxed{\mathrm{c}_i} & \boxed{\mathrm{v}_t} & \boxed{\mathrm{c}_o}\\
            \mathrm{\varnothing} & \mathrm{wa} & \underbrace{\mathrm{dd}}_{\text{compressed}}\\
            \mathrm{} & \mathrm{} & \mathrm{}
        \end{pmatrix}
        =
        \begin{pmatrix}\pmb{\text{\Large 왔}}\\\boxed{\mathrm{syllable}}\\\mathrm{wadd}\\\mathrm{}\end{pmatrix}
    \end{align}\\
    
    \begin{align}
        \textbf{\Huge 닭} &= 
        \underbrace{
            \begin{pmatrix}\pmb{\text{\Large ㄷ}}\\\boxed{\mathrm{c}_1}\\\mathrm{d}\\\mathrm{}\end{pmatrix}
        }_{\mathrm{c}_i}
        +
        \underbrace{
            \begin{pmatrix}             
                \pmb{\text{\Large }} & \pmb{\text{\Large ㅏ}}\\
                \boxed{\mathrm{v}_1} & \boxed{\mathrm{v}_2}\\
                \mathrm{} & \mathrm{a}\\
                \mathrm{} & \mathrm{}
                \end{pmatrix}
        }_{\mathrm{v}_t}
            +
        \underbrace{
            \begin{pmatrix}
                \pmb{\text{\Large ㄹ}} & \pmb{\text{\Large ㄱ}}\\
                \boxed{\mathrm{c}_2} & \boxed{\mathrm{c}_3}\\
                \mathrm{l} & \mathrm{g}\\
                \mathrm{} & \mathrm{}
            \end{pmatrix}
        }_{\mathrm{c}_o}
        \\&=
        \begin{pmatrix}
            \pmb{\text{\Large ㄷ}} & \pmb{\text{\Large ㅏ}} & \pmb{\text{\Large ㄺ}}\\
            \boxed{\mathrm{c}_i} & \boxed{\mathrm{v}_t} & \boxed{\mathrm{c}_o}\\
            \mathrm{d} & \mathrm{a} & \underbrace{\mathrm{lg}}_{\text{compressed}}\\
            \mathrm{} & \mathrm{} & \mathrm{}
        \end{pmatrix}
        =
        \begin{pmatrix}\pmb{\text{\Large 닭}}\\\boxed{\mathrm{syllable}}\\\mathrm{dalg}\\\mathrm{}\end{pmatrix}
    \end{align}
    \subsection{With special cases}
    \begin{align}
        \textbf{\Huge 시} &\overset{\text{special case}}{\underset{\text{\cref{eq:shieut}}}{=}} 
        \underbrace{
            \begin{pmatrix}\pmb{\text{\Large ㅅ}}\\\boxed{\mathrm{c}_1}\\\mathrm{s}\\\mathrm{sh}\end{pmatrix}
        }_{\mathrm{c}_i}
        +
        \underbrace{
            \begin{pmatrix}             
                \pmb{\text{\Large }} & \pmb{\text{\Large ㅣ}}\\
                \boxed{\mathrm{v}_1} & \boxed{\mathrm{v}_2}\\
                \mathrm{} & \mathrm{i}\\
                \mathrm{} & \mathrm{}
                \end{pmatrix}
        }_{\mathrm{v}_t}
            +
        \underbrace{
            \begin{pmatrix}
                \pmb{\text{\Large }} & \pmb{\text{\Large }}\\
                \boxed{\mathrm{c}_2} & \boxed{\mathrm{c}_3}\\
                \mathrm{} & \mathrm{}\\
                \mathrm{} & \mathrm{}
            \end{pmatrix}
        }_{\mathrm{c}_o}
        \\&=
        \begin{pmatrix}
            \pmb{\text{\Large ㅅ}} & \pmb{\text{\Large ㅣ}} & \pmb{\text{\Large }}\\
            \boxed{\mathrm{c}_i} & \boxed{\mathrm{v}_t} & \boxed{\mathrm{c}_o}\\
            \mathrm{s} & \mathrm{i} & \mathrm{}\\
            \mathrm{sh} & \mathrm{} & \mathrm{}
        \end{pmatrix}
        =
        \begin{pmatrix}\pmb{\text{\Large 시}}\\\boxed{\mathrm{syllable}}\\\mathrm{si}\ \pmb{\times}\\\mathrm{shi}\ \checkmark\end{pmatrix}
    \end{align}\\    
    
    \begin{align}
        \underbrace{\text{\Huge 빨}}_{\text{bbal}}\underbrace{\textbf{\Huge 리}}_{\text{li}} &\overset{\text{special case}}{\underset{\text{\cref{eq:rieut}}}{=}} 
        \underbrace{
            \begin{pmatrix}\pmb{\text{\Large ㄹ}}\\\boxed{\mathrm{c}_1}\\\mathrm{r}\\\mathrm{l}\end{pmatrix}
        }_{\mathrm{c}_i}
        +
        \underbrace{
            \begin{pmatrix}             
                \pmb{\text{\Large }} & \pmb{\text{\Large ㅣ}}\\
                \boxed{\mathrm{v}_1} & \boxed{\mathrm{v}_2}\\
                \mathrm{} & \mathrm{i}\\
                \mathrm{} & \mathrm{}
                \end{pmatrix}
        }_{\mathrm{v}_t}
            +
        \underbrace{
            \begin{pmatrix}
                \pmb{\text{\Large }} & \pmb{\text{\Large }}\\
                \boxed{\mathrm{c}_2} & \boxed{\mathrm{c}_3}\\
                \mathrm{} & \mathrm{}\\
                \mathrm{} & \mathrm{}
            \end{pmatrix}
        }_{\mathrm{c}_o}
        \\&=
        \begin{pmatrix}
            \pmb{\text{\Large ㄹ}} & \pmb{\text{\Large ㅣ}} & \pmb{\text{\Large }}\\
            \boxed{\mathrm{c}_i} & \boxed{\mathrm{v}_t} & \boxed{\mathrm{c}_o}\\
            \mathrm{r} & \mathrm{i} & \mathrm{}\\
            \mathrm{l} & \mathrm{} & \mathrm{}
        \end{pmatrix}
        =
        \begin{pmatrix}\pmb{\text{\Large 리}}\\\boxed{\mathrm{syllable}}\\\mathrm{ri}\ \pmb{\times}\\\mathrm{li}\ \checkmark \end{pmatrix}
    \end{align}\\
    
    \begin{align}
        \underbrace{\textbf{\Huge 있}}_{\text{iss}}\underbrace{\text{\Huge 어}}_{\text{uh}} &\overset{\text{special case}}{\underset{\text{\cref{eq:shieut}}}{=}}  
        \underbrace{
            \begin{pmatrix}\pmb{\text{\Large ㅇ}}\\\boxed{\mathrm{c}_1}\\\mathrm{\varnothing}\\\mathrm{}\end{pmatrix}
        }_{\mathrm{c}_i}
        +
        \underbrace{
            \begin{pmatrix}             
                \pmb{\text{\Large }} & \pmb{\text{\Large ㅣ}}\\
                \boxed{\mathrm{v}_1} & \boxed{\mathrm{v}_2}\\
                \mathrm{} & \mathrm{i}\\
                \mathrm{} & \mathrm{}
                \end{pmatrix}
        }_{\mathrm{v}_t}
            +
        \underbrace{
            \begin{pmatrix}
                \pmb{\text{\Large ㅅ}} & \pmb{\text{\Large ㅅ}}\\
                \boxed{\mathrm{c}_2} & \boxed{\mathrm{c}_3}\\
                \mathrm{d} & \mathrm{d}\\
                \mathrm{s} & \mathrm{s}
            \end{pmatrix}
        }_{\mathrm{c}_o}
        \\&=
        \begin{pmatrix}
            \pmb{\text{\Large ㅇ}} & \pmb{\text{\Large ㅣ}} & \pmb{\text{\Large ㅆ}}\\
            \boxed{\mathrm{c}_i} & \boxed{\mathrm{v}_t} & \boxed{\mathrm{c}_o}\\
            \mathrm{\varnothing} & \mathrm{i} & \underbrace{\mathrm{dd}}_{\text{compressed}}\\
            \mathrm{} & \mathrm{} & \underbrace{\mathrm{ss}}_{\text{compressed}}
        \end{pmatrix}
        =
        \begin{pmatrix}\pmb{\text{\Large 있}}\\\boxed{\mathrm{syllable}}\\\mathrm{idd}\ \pmb{\times}\\\mathrm{iss}\ \checkmark \end{pmatrix}
    \end{align}

    
    
    
\section{Exercises}
    Reminder.
    \begin{figure}[H]
        \centering
        \begin{tikzpicture}[scale=0.9]
            \draw (0,2) rectangle (2,3) node[pos=.5] {$c_1$};
            \draw (0,1) rectangle (2,2) node[pos=.5] {$v_1$};
            \draw (2,1) rectangle (3,3) node[pos=.5] {$v_2$};
            \draw (0,0) rectangle (1.5,1) node[pos=.5] {$c_2$};
            \draw (1.5,0) rectangle (3,1) node[pos=.5] {$c_3$};
            \node[anchor=west, align=left] at (3.4,1.5) {$c$\ldots consonant\\ $v$\ldots vowel};
        \end{tikzpicture}
        \caption*{Reminder: Syllable Block}
    \end{figure}

    \begin{equation}
        \begin{pmatrix} \pmb{\text{\Large hangul letter}} \\ \boxed{\text{position in Block}} \\ \text{default sound} \\ \text{special case}\end{pmatrix}
    \end{equation}

    \begin{align}
        \textbf{\Huge } &= 
        \underbrace{
            \begin{pmatrix}\pmb{\text{\Large }}\\\boxed{\mathrm{c}_1}\\ \mathrm{}\\ \mathrm{} \end{pmatrix}
        }_{\mathrm{c}_i}
        +
        \underbrace{
            \begin{pmatrix}             
                \pmb{\text{\Large }}&\pmb{\text{\Large }}\\
                \boxed{\mathrm{v}_1}&\boxed{\mathrm{v}_2}\\
                \mathrm{}&\mathrm{}\\
                \mathrm{}&\mathrm{}
            \end{pmatrix}
        }_{\mathrm{v}_t}
        +
        \underbrace{
            \begin{pmatrix}
                \pmb{\text{\Large }} & \pmb{\text{\Large }}\\
                \boxed{\mathrm{c}_2} & \boxed{\mathrm{c}_3}\\
                \mathrm{} & \mathrm{}\\
                \mathrm{} & \mathrm{}
            \end{pmatrix}
        }_{\mathrm{c}_o}
        \\&=
        \begin{pmatrix}
            \pmb{\text{\Large }} & \pmb{\text{\Large }} & \pmb{\text{\Large }}\\
            \boxed{\mathrm{c}_i} & \boxed{\mathrm{v}_t} & \boxed{\mathrm{c}_o}\\
            \mathrm{} & \mathrm{} & \mathrm{}\\
            \mathrm{} & \mathrm{} & \mathrm{}
        \end{pmatrix}
        =
        \begin{pmatrix}\pmb{\text{\Large }}\\\boxed{\mathrm{syllable}}\\\mathrm{}\\\mathrm{}\end{pmatrix}_{\text{\footnotesize Answer:}}
    \end{align}
    
    \subsection{Initial consonant + single vowel}
    \begin{align}
        \textbf{\Huge 나} &= 
        \underbrace{
            \begin{pmatrix}\pmb{\text{\Large }}\\\boxed{\mathrm{c}_1}\\ \mathrm{}\\ \mathrm{} \end{pmatrix}
        }_{\mathrm{c}_i}
        +
        \underbrace{
            \begin{pmatrix}             
                \pmb{\text{\Large }}&\pmb{\text{\Large }}\\
                \boxed{\mathrm{v}_1}&\boxed{\mathrm{v}_2}\\
                \mathrm{}&\mathrm{}\\
                \mathrm{}&\mathrm{}
            \end{pmatrix}
        }_{\mathrm{v}_t}
        +
        \underbrace{
            \begin{pmatrix}
                \pmb{\text{\Large }} & \pmb{\text{\Large }}\\
                \boxed{\mathrm{c}_2} & \boxed{\mathrm{c}_3}\\
                \mathrm{} & \mathrm{}\\
                \mathrm{} & \mathrm{}
            \end{pmatrix}
        }_{\mathrm{c}_o}
        \\&=
        \begin{pmatrix}
            \pmb{\text{\Large }} & \pmb{\text{\Large }} & \pmb{\text{\Large }}\\
            \boxed{\mathrm{c}_i} & \boxed{\mathrm{v}_t} & \boxed{\mathrm{c}_o}\\
            \mathrm{} & \mathrm{} & \mathrm{}\\
            \mathrm{} & \mathrm{} & \mathrm{}
        \end{pmatrix}
        =
        \begin{pmatrix}\pmb{\text{\Large }}\\\boxed{\mathrm{syllable}}\\\mathrm{}\\\mathrm{}\end{pmatrix}_{\text{\footnotesize Answer: na}}
    \end{align}\\

    \begin{align}
        \textbf{\Huge 미} &= 
        \underbrace{
            \begin{pmatrix}\pmb{\text{\Large }}\\\boxed{\mathrm{c}_1}\\ \mathrm{}\\ \mathrm{} \end{pmatrix}
        }_{\mathrm{c}_i}
        +
        \underbrace{
            \begin{pmatrix}             
                \pmb{\text{\Large }}&\pmb{\text{\Large }}\\
                \boxed{\mathrm{v}_1}&\boxed{\mathrm{v}_2}\\
                \mathrm{}&\mathrm{}\\
                \mathrm{}&\mathrm{}
            \end{pmatrix}
        }_{\mathrm{v}_t}
        +
        \underbrace{
            \begin{pmatrix}
                \pmb{\text{\Large }} & \pmb{\text{\Large }}\\
                \boxed{\mathrm{c}_2} & \boxed{\mathrm{c}_3}\\
                \mathrm{} & \mathrm{}\\
                \mathrm{} & \mathrm{}
            \end{pmatrix}
        }_{\mathrm{c}_o}
        \\&=
        \begin{pmatrix}
            \pmb{\text{\Large }} & \pmb{\text{\Large }} & \pmb{\text{\Large }}\\
            \boxed{\mathrm{c}_i} & \boxed{\mathrm{v}_t} & \boxed{\mathrm{c}_o}\\
            \mathrm{} & \mathrm{} & \mathrm{}\\
            \mathrm{} & \mathrm{} & \mathrm{}
        \end{pmatrix}
        =
        \begin{pmatrix}\pmb{\text{\Large }}\\\boxed{\mathrm{syllable}}\\\mathrm{}\\\mathrm{}\end{pmatrix}_{\text{\footnotesize Answer: mi}}
    \end{align}\\

    \begin{align}
        \textbf{\Huge 개} &= 
        \underbrace{
            \begin{pmatrix}\pmb{\text{\Large }}\\\boxed{\mathrm{c}_1}\\ \mathrm{}\\ \mathrm{} \end{pmatrix}
        }_{\mathrm{c}_i}
        +
        \underbrace{
            \begin{pmatrix}             
                \pmb{\text{\Large }}&\pmb{\text{\Large }}\\
                \boxed{\mathrm{v}_1}&\boxed{\mathrm{v}_2}\\
                \mathrm{}&\mathrm{}\\
                \mathrm{}&\mathrm{}
            \end{pmatrix}
        }_{\mathrm{v}_t}
        +
        \underbrace{
            \begin{pmatrix}
                \pmb{\text{\Large }} & \pmb{\text{\Large }}\\
                \boxed{\mathrm{c}_2} & \boxed{\mathrm{c}_3}\\
                \mathrm{} & \mathrm{}\\
                \mathrm{} & \mathrm{}
            \end{pmatrix}
        }_{\mathrm{c}_o}
        \\&=
        \begin{pmatrix}
            \pmb{\text{\Large }} & \pmb{\text{\Large }} & \pmb{\text{\Large }}\\
            \boxed{\mathrm{c}_i} & \boxed{\mathrm{v}_t} & \boxed{\mathrm{c}_o}\\
            \mathrm{} & \mathrm{} & \mathrm{}\\
            \mathrm{} & \mathrm{} & \mathrm{}
        \end{pmatrix}
        =
        \begin{pmatrix}\pmb{\text{\Large }}\\\boxed{\mathrm{syllable}}\\\mathrm{}\\\mathrm{}\end{pmatrix}_{\text{\footnotesize Answer: gae}}
    \end{align}\\

    \begin{align}
        \textbf{\Huge 너} &= 
        \underbrace{
            \begin{pmatrix}\pmb{\text{\Large }}\\\boxed{\mathrm{c}_1}\\ \mathrm{}\\ \mathrm{} \end{pmatrix}
        }_{\mathrm{c}_i}
        +
        \underbrace{
            \begin{pmatrix}             
                \pmb{\text{\Large }}&\pmb{\text{\Large }}\\
                \boxed{\mathrm{v}_1}&\boxed{\mathrm{v}_2}\\
                \mathrm{}&\mathrm{}\\
                \mathrm{}&\mathrm{}
            \end{pmatrix}
        }_{\mathrm{v}_t}
        +
        \underbrace{
            \begin{pmatrix}
                \pmb{\text{\Large }} & \pmb{\text{\Large }}\\
                \boxed{\mathrm{c}_2} & \boxed{\mathrm{c}_3}\\
                \mathrm{} & \mathrm{}\\
                \mathrm{} & \mathrm{}
            \end{pmatrix}
        }_{\mathrm{c}_o}
        \\&=
        \begin{pmatrix}
            \pmb{\text{\Large }} & \pmb{\text{\Large }} & \pmb{\text{\Large }}\\
            \boxed{\mathrm{c}_i} & \boxed{\mathrm{v}_t} & \boxed{\mathrm{c}_o}\\
            \mathrm{} & \mathrm{} & \mathrm{}\\
            \mathrm{} & \mathrm{} & \mathrm{}
        \end{pmatrix}
        =
        \begin{pmatrix}\pmb{\text{\Large }}\\\boxed{\mathrm{syllable}}\\\mathrm{}\\\mathrm{}\end{pmatrix}_{\text{\footnotesize Answer: nuh}}
    \end{align}\\

    \begin{align}
        \textbf{\Huge 도} &= 
        \underbrace{
            \begin{pmatrix}\pmb{\text{\Large }}\\\boxed{\mathrm{c}_1}\\ \mathrm{}\\ \mathrm{} \end{pmatrix}
        }_{\mathrm{c}_i}
        +
        \underbrace{
            \begin{pmatrix}             
                \pmb{\text{\Large }}&\pmb{\text{\Large }}\\
                \boxed{\mathrm{v}_1}&\boxed{\mathrm{v}_2}\\
                \mathrm{}&\mathrm{}\\
                \mathrm{}&\mathrm{}
            \end{pmatrix}
        }_{\mathrm{v}_t}
        +
        \underbrace{
            \begin{pmatrix}
                \pmb{\text{\Large }} & \pmb{\text{\Large }}\\
                \boxed{\mathrm{c}_2} & \boxed{\mathrm{c}_3}\\
                \mathrm{} & \mathrm{}\\
                \mathrm{} & \mathrm{}
            \end{pmatrix}
        }_{\mathrm{c}_o}
        \\&=
        \begin{pmatrix}
            \pmb{\text{\Large }} & \pmb{\text{\Large }} & \pmb{\text{\Large }}\\
            \boxed{\mathrm{c}_i} & \boxed{\mathrm{v}_t} & \boxed{\mathrm{c}_o}\\
            \mathrm{} & \mathrm{} & \mathrm{}\\
            \mathrm{} & \mathrm{} & \mathrm{}
        \end{pmatrix}
        =
        \begin{pmatrix}\pmb{\text{\Large }}\\\boxed{\mathrm{syllable}}\\\mathrm{}\\\mathrm{}\end{pmatrix}_{\text{\footnotesize Answer: do}}
    \end{align}\\

    \begin{align}
        \textbf{\Huge 그} &= 
        \underbrace{
            \begin{pmatrix}\pmb{\text{\Large }}\\\boxed{\mathrm{c}_1}\\ \mathrm{}\\ \mathrm{} \end{pmatrix}
        }_{\mathrm{c}_i}
        +
        \underbrace{
            \begin{pmatrix}             
                \pmb{\text{\Large }}&\pmb{\text{\Large }}\\
                \boxed{\mathrm{v}_1}&\boxed{\mathrm{v}_2}\\
                \mathrm{}&\mathrm{}\\
                \mathrm{}&\mathrm{}
            \end{pmatrix}
        }_{\mathrm{v}_t}
        +
        \underbrace{
            \begin{pmatrix}
                \pmb{\text{\Large }} & \pmb{\text{\Large }}\\
                \boxed{\mathrm{c}_2} & \boxed{\mathrm{c}_3}\\
                \mathrm{} & \mathrm{}\\
                \mathrm{} & \mathrm{}
            \end{pmatrix}
        }_{\mathrm{c}_o}
        \\&=
        \begin{pmatrix}
            \pmb{\text{\Large }} & \pmb{\text{\Large }} & \pmb{\text{\Large }}\\
            \boxed{\mathrm{c}_i} & \boxed{\mathrm{v}_t} & \boxed{\mathrm{c}_o}\\
            \mathrm{} & \mathrm{} & \mathrm{}\\
            \mathrm{} & \mathrm{} & \mathrm{}
        \end{pmatrix}
        =
        \begin{pmatrix}\pmb{\text{\Large }}\\\boxed{\mathrm{syllable}}\\\mathrm{}\\\mathrm{}\end{pmatrix}_{\text{\footnotesize Answer: geww}}
    \end{align}\\

    \begin{align}
        \textbf{\Huge 효} &= 
        \underbrace{
            \begin{pmatrix}\pmb{\text{\Large }}\\\boxed{\mathrm{c}_1}\\ \mathrm{}\\ \mathrm{} \end{pmatrix}
        }_{\mathrm{c}_i}
        +
        \underbrace{
            \begin{pmatrix}             
                \pmb{\text{\Large }}&\pmb{\text{\Large }}\\
                \boxed{\mathrm{v}_1}&\boxed{\mathrm{v}_2}\\
                \mathrm{}&\mathrm{}\\
                \mathrm{}&\mathrm{}
            \end{pmatrix}
        }_{\mathrm{v}_t}
        +
        \underbrace{
            \begin{pmatrix}
                \pmb{\text{\Large }} & \pmb{\text{\Large }}\\
                \boxed{\mathrm{c}_2} & \boxed{\mathrm{c}_3}\\
                \mathrm{} & \mathrm{}\\
                \mathrm{} & \mathrm{}
            \end{pmatrix}
        }_{\mathrm{c}_o}
        \\&=
        \begin{pmatrix}
            \pmb{\text{\Large }} & \pmb{\text{\Large }} & \pmb{\text{\Large }}\\
            \boxed{\mathrm{c}_i} & \boxed{\mathrm{v}_t} & \boxed{\mathrm{c}_o}\\
            \mathrm{} & \mathrm{} & \mathrm{}\\
            \mathrm{} & \mathrm{} & \mathrm{}
        \end{pmatrix}
        =
        \begin{pmatrix}\pmb{\text{\Large }}\\\boxed{\mathrm{syllable}}\\\mathrm{}\\\mathrm{}\end{pmatrix}_{\text{\footnotesize Answer: hyo}}
    \end{align}
    \subsection{Initial consonant + compound vowel}
    \begin{align}
        \textbf{\Huge 왜} &= 
        \underbrace{
            \begin{pmatrix}\pmb{\text{\Large }}\\\boxed{\mathrm{c}_1}\\ \mathrm{}\\ \mathrm{} \end{pmatrix}
        }_{\mathrm{c}_i}
        +
        \underbrace{
            \begin{pmatrix}             
                \pmb{\text{\Large }}&\pmb{\text{\Large }}\\
                \boxed{\mathrm{v}_1}&\boxed{\mathrm{v}_2}\\
                \mathrm{}&\mathrm{}\\
                \mathrm{}&\mathrm{}
            \end{pmatrix}
        }_{\mathrm{v}_t}
        +
        \underbrace{
            \begin{pmatrix}
                \pmb{\text{\Large }} & \pmb{\text{\Large }}\\
                \boxed{\mathrm{c}_2} & \boxed{\mathrm{c}_3}\\
                \mathrm{} & \mathrm{}\\
                \mathrm{} & \mathrm{}
            \end{pmatrix}
        }_{\mathrm{c}_o}
        \\&=
        \begin{pmatrix}
            \pmb{\text{\Large }} & \pmb{\text{\Large }} & \pmb{\text{\Large }}\\
            \boxed{\mathrm{c}_i} & \boxed{\mathrm{v}_t} & \boxed{\mathrm{c}_o}\\
            \mathrm{} & \mathrm{} & \mathrm{}\\
            \mathrm{} & \mathrm{} & \mathrm{}
        \end{pmatrix}
        =
        \begin{pmatrix}\pmb{\text{\Large }}\\\boxed{\mathrm{syllable}}\\\mathrm{}\\\mathrm{}\end{pmatrix}_{\text{\footnotesize Answer: wae}}
    \end{align}\\

    \begin{align}
        \textbf{\Huge 위} &= 
        \underbrace{
            \begin{pmatrix}\pmb{\text{\Large }}\\\boxed{\mathrm{c}_1}\\ \mathrm{}\\ \mathrm{} \end{pmatrix}
        }_{\mathrm{c}_i}
        +
        \underbrace{
            \begin{pmatrix}             
                \pmb{\text{\Large }}&\pmb{\text{\Large }}\\
                \boxed{\mathrm{v}_1}&\boxed{\mathrm{v}_2}\\
                \mathrm{}&\mathrm{}\\
                \mathrm{}&\mathrm{}
            \end{pmatrix}
        }_{\mathrm{v}_t}
        +
        \underbrace{
            \begin{pmatrix}
                \pmb{\text{\Large }} & \pmb{\text{\Large }}\\
                \boxed{\mathrm{c}_2} & \boxed{\mathrm{c}_3}\\
                \mathrm{} & \mathrm{}\\
                \mathrm{} & \mathrm{}
            \end{pmatrix}
        }_{\mathrm{c}_o}
        \\&=
        \begin{pmatrix}
            \pmb{\text{\Large }} & \pmb{\text{\Large }} & \pmb{\text{\Large }}\\
            \boxed{\mathrm{c}_i} & \boxed{\mathrm{v}_t} & \boxed{\mathrm{c}_o}\\
            \mathrm{} & \mathrm{} & \mathrm{}\\
            \mathrm{} & \mathrm{} & \mathrm{}
        \end{pmatrix}
        =
        \begin{pmatrix}\pmb{\text{\Large }}\\\boxed{\mathrm{syllable}}\\\mathrm{}\\\mathrm{}\end{pmatrix}_{\text{\footnotesize Answer: wi}}
    \end{align}\\

    \begin{align}
        \textbf{\Huge 뭐} &= 
        \underbrace{
            \begin{pmatrix}\pmb{\text{\Large }}\\\boxed{\mathrm{c}_1}\\ \mathrm{}\\ \mathrm{} \end{pmatrix}
        }_{\mathrm{c}_i}
        +
        \underbrace{
            \begin{pmatrix}             
                \pmb{\text{\Large }}&\pmb{\text{\Large }}\\
                \boxed{\mathrm{v}_1}&\boxed{\mathrm{v}_2}\\
                \mathrm{}&\mathrm{}\\
                \mathrm{}&\mathrm{}
            \end{pmatrix}
        }_{\mathrm{v}_t}
        +
        \underbrace{
            \begin{pmatrix}
                \pmb{\text{\Large }} & \pmb{\text{\Large }}\\
                \boxed{\mathrm{c}_2} & \boxed{\mathrm{c}_3}\\
                \mathrm{} & \mathrm{}\\
                \mathrm{} & \mathrm{}
            \end{pmatrix}
        }_{\mathrm{c}_o}
        \\&=
        \begin{pmatrix}
            \pmb{\text{\Large }} & \pmb{\text{\Large }} & \pmb{\text{\Large }}\\
            \boxed{\mathrm{c}_i} & \boxed{\mathrm{v}_t} & \boxed{\mathrm{c}_o}\\
            \mathrm{} & \mathrm{} & \mathrm{}\\
            \mathrm{} & \mathrm{} & \mathrm{}
        \end{pmatrix}
        =
        \begin{pmatrix}\pmb{\text{\Large }}\\\boxed{\mathrm{syllable}}\\\mathrm{}\\\mathrm{}\end{pmatrix}_{\text{\footnotesize Answer: mwo}}
    \end{align}\\

    \begin{align}
        \textbf{\Huge 의} &= 
        \underbrace{
            \begin{pmatrix}\pmb{\text{\Large }}\\\boxed{\mathrm{c}_1}\\ \mathrm{}\\ \mathrm{} \end{pmatrix}
        }_{\mathrm{c}_i}
        +
        \underbrace{
            \begin{pmatrix}             
                \pmb{\text{\Large }}&\pmb{\text{\Large }}\\
                \boxed{\mathrm{v}_1}&\boxed{\mathrm{v}_2}\\
                \mathrm{}&\mathrm{}\\
                \mathrm{}&\mathrm{}
            \end{pmatrix}
        }_{\mathrm{v}_t}
        +
        \underbrace{
            \begin{pmatrix}
                \pmb{\text{\Large }} & \pmb{\text{\Large }}\\
                \boxed{\mathrm{c}_2} & \boxed{\mathrm{c}_3}\\
                \mathrm{} & \mathrm{}\\
                \mathrm{} & \mathrm{}
            \end{pmatrix}
        }_{\mathrm{c}_o}
        \\&=
        \begin{pmatrix}
            \pmb{\text{\Large }} & \pmb{\text{\Large }} & \pmb{\text{\Large }}\\
            \boxed{\mathrm{c}_i} & \boxed{\mathrm{v}_t} & \boxed{\mathrm{c}_o}\\
            \mathrm{} & \mathrm{} & \mathrm{}\\
            \mathrm{} & \mathrm{} & \mathrm{}
        \end{pmatrix}
        =
        \begin{pmatrix}\pmb{\text{\Large }}\\\boxed{\mathrm{syllable}}\\\mathrm{}\\\mathrm{}\end{pmatrix}_{\text{\footnotesize Answer: ui}}
    \end{align}
    \subsection{Initial consonant + single vowel + optional consonant}
    \begin{align}
        \textbf{\Huge 산} &= 
        \underbrace{
            \begin{pmatrix}\pmb{\text{\Large }}\\\boxed{\mathrm{c}_1}\\ \mathrm{}\\ \mathrm{} \end{pmatrix}
        }_{\mathrm{c}_i}
        +
        \underbrace{
            \begin{pmatrix}             
                \pmb{\text{\Large }}&\pmb{\text{\Large }}\\
                \boxed{\mathrm{v}_1}&\boxed{\mathrm{v}_2}\\
                \mathrm{}&\mathrm{}\\
                \mathrm{}&\mathrm{}
            \end{pmatrix}
        }_{\mathrm{v}_t}
        +
        \underbrace{
            \begin{pmatrix}
                \pmb{\text{\Large }} & \pmb{\text{\Large }}\\
                \boxed{\mathrm{c}_2} & \boxed{\mathrm{c}_3}\\
                \mathrm{} & \mathrm{}\\
                \mathrm{} & \mathrm{}
            \end{pmatrix}
        }_{\mathrm{c}_o}
        \\&=
        \begin{pmatrix}
            \pmb{\text{\Large }} & \pmb{\text{\Large }} & \pmb{\text{\Large }}\\
            \boxed{\mathrm{c}_i} & \boxed{\mathrm{v}_t} & \boxed{\mathrm{c}_o}\\
            \mathrm{} & \mathrm{} & \mathrm{}\\
            \mathrm{} & \mathrm{} & \mathrm{}
        \end{pmatrix}
        =
        \begin{pmatrix}\pmb{\text{\Large }}\\\boxed{\mathrm{syllable}}\\\mathrm{}\\\mathrm{}\end{pmatrix}_{\text{\footnotesize Answer: san}}
    \end{align}\\

    \begin{align}
        \textbf{\Huge 물} &= 
        \underbrace{
            \begin{pmatrix}\pmb{\text{\Large }}\\\boxed{\mathrm{c}_1}\\ \mathrm{}\\ \mathrm{} \end{pmatrix}
        }_{\mathrm{c}_i}
        +
        \underbrace{
            \begin{pmatrix}             
                \pmb{\text{\Large }}&\pmb{\text{\Large }}\\
                \boxed{\mathrm{v}_1}&\boxed{\mathrm{v}_2}\\
                \mathrm{}&\mathrm{}\\
                \mathrm{}&\mathrm{}
            \end{pmatrix}
        }_{\mathrm{v}_t}
        +
        \underbrace{
            \begin{pmatrix}
                \pmb{\text{\Large }} & \pmb{\text{\Large }}\\
                \boxed{\mathrm{c}_2} & \boxed{\mathrm{c}_3}\\
                \mathrm{} & \mathrm{}\\
                \mathrm{} & \mathrm{}
            \end{pmatrix}
        }_{\mathrm{c}_o}
        \\&=
        \begin{pmatrix}
            \pmb{\text{\Large }} & \pmb{\text{\Large }} & \pmb{\text{\Large }}\\
            \boxed{\mathrm{c}_i} & \boxed{\mathrm{v}_t} & \boxed{\mathrm{c}_o}\\
            \mathrm{} & \mathrm{} & \mathrm{}\\
            \mathrm{} & \mathrm{} & \mathrm{}
        \end{pmatrix}
        =
        \begin{pmatrix}\pmb{\text{\Large }}\\\boxed{\mathrm{syllable}}\\\mathrm{}\\\mathrm{}\end{pmatrix}_{\text{\footnotesize Answer: mool}}
    \end{align}\\

    \begin{align}
        \textbf{\Huge 강} &= 
        \underbrace{
            \begin{pmatrix}\pmb{\text{\Large }}\\\boxed{\mathrm{c}_1}\\ \mathrm{}\\ \mathrm{} \end{pmatrix}
        }_{\mathrm{c}_i}
        +
        \underbrace{
            \begin{pmatrix}             
                \pmb{\text{\Large }}&\pmb{\text{\Large }}\\
                \boxed{\mathrm{v}_1}&\boxed{\mathrm{v}_2}\\
                \mathrm{}&\mathrm{}\\
                \mathrm{}&\mathrm{}
            \end{pmatrix}
        }_{\mathrm{v}_t}
        +
        \underbrace{
            \begin{pmatrix}
                \pmb{\text{\Large }} & \pmb{\text{\Large }}\\
                \boxed{\mathrm{c}_2} & \boxed{\mathrm{c}_3}\\
                \mathrm{} & \mathrm{}\\
                \mathrm{} & \mathrm{}
            \end{pmatrix}
        }_{\mathrm{c}_o}
        \\&=
        \begin{pmatrix}
            \pmb{\text{\Large }} & \pmb{\text{\Large }} & \pmb{\text{\Large }}\\
            \boxed{\mathrm{c}_i} & \boxed{\mathrm{v}_t} & \boxed{\mathrm{c}_o}\\
            \mathrm{} & \mathrm{} & \mathrm{}\\
            \mathrm{} & \mathrm{} & \mathrm{}
        \end{pmatrix}
        =
        \begin{pmatrix}\pmb{\text{\Large }}\\\boxed{\mathrm{syllable}}\\\mathrm{}\\\mathrm{}\end{pmatrix}_{\text{\footnotesize Answer: gang}}
    \end{align}\\

    \begin{align}
        \textbf{\Huge 옷} &= 
        \underbrace{
            \begin{pmatrix}\pmb{\text{\Large }}\\\boxed{\mathrm{c}_1}\\ \mathrm{}\\ \mathrm{} \end{pmatrix}
        }_{\mathrm{c}_i}
        +
        \underbrace{
            \begin{pmatrix}             
                \pmb{\text{\Large }}&\pmb{\text{\Large }}\\
                \boxed{\mathrm{v}_1}&\boxed{\mathrm{v}_2}\\
                \mathrm{}&\mathrm{}\\
                \mathrm{}&\mathrm{}
            \end{pmatrix}
        }_{\mathrm{v}_t}
        +
        \underbrace{
            \begin{pmatrix}
                \pmb{\text{\Large }} & \pmb{\text{\Large }}\\
                \boxed{\mathrm{c}_2} & \boxed{\mathrm{c}_3}\\
                \mathrm{} & \mathrm{}\\
                \mathrm{} & \mathrm{}
            \end{pmatrix}
        }_{\mathrm{c}_o}
        \\&=
        \begin{pmatrix}
            \pmb{\text{\Large }} & \pmb{\text{\Large }} & \pmb{\text{\Large }}\\
            \boxed{\mathrm{c}_i} & \boxed{\mathrm{v}_t} & \boxed{\mathrm{c}_o}\\
            \mathrm{} & \mathrm{} & \mathrm{}\\
            \mathrm{} & \mathrm{} & \mathrm{}
        \end{pmatrix}
        =
        \begin{pmatrix}\pmb{\text{\Large }}\\\boxed{\mathrm{syllable}}\\\mathrm{}\\\mathrm{}\end{pmatrix}_{\text{\footnotesize Answer: od}}
    \end{align}\\

    \begin{align}
        \textbf{\Huge 낮} &= 
        \underbrace{
            \begin{pmatrix}\pmb{\text{\Large }}\\\boxed{\mathrm{c}_1}\\ \mathrm{}\\ \mathrm{} \end{pmatrix}
        }_{\mathrm{c}_i}
        +
        \underbrace{
            \begin{pmatrix}             
                \pmb{\text{\Large }}&\pmb{\text{\Large }}\\
                \boxed{\mathrm{v}_1}&\boxed{\mathrm{v}_2}\\
                \mathrm{}&\mathrm{}\\
                \mathrm{}&\mathrm{}
            \end{pmatrix}
        }_{\mathrm{v}_t}
        +
        \underbrace{
            \begin{pmatrix}
                \pmb{\text{\Large }} & \pmb{\text{\Large }}\\
                \boxed{\mathrm{c}_2} & \boxed{\mathrm{c}_3}\\
                \mathrm{} & \mathrm{}\\
                \mathrm{} & \mathrm{}
            \end{pmatrix}
        }_{\mathrm{c}_o}
        \\&=
        \begin{pmatrix}
            \pmb{\text{\Large }} & \pmb{\text{\Large }} & \pmb{\text{\Large }}\\
            \boxed{\mathrm{c}_i} & \boxed{\mathrm{v}_t} & \boxed{\mathrm{c}_o}\\
            \mathrm{} & \mathrm{} & \mathrm{}\\
            \mathrm{} & \mathrm{} & \mathrm{}
        \end{pmatrix}
        =
        \begin{pmatrix}\pmb{\text{\Large }}\\\boxed{\mathrm{syllable}}\\\mathrm{}\\\mathrm{}\end{pmatrix}_{\text{\footnotesize Answer: nad}}
    \end{align}
    
    \subsection{Initial consonant + compound vowel + optional consonant}
    \begin{align}
        \textbf{\Huge 원} &= 
        \underbrace{
            \begin{pmatrix}\pmb{\text{\Large }}\\\boxed{\mathrm{c}_1}\\ \mathrm{}\\ \mathrm{} \end{pmatrix}
        }_{\mathrm{c}_i}
        +
        \underbrace{
            \begin{pmatrix}             
                \pmb{\text{\Large }}&\pmb{\text{\Large }}\\
                \boxed{\mathrm{v}_1}&\boxed{\mathrm{v}_2}\\
                \mathrm{}&\mathrm{}\\
                \mathrm{}&\mathrm{}
            \end{pmatrix}
        }_{\mathrm{v}_t}
        +
        \underbrace{
            \begin{pmatrix}
                \pmb{\text{\Large }} & \pmb{\text{\Large }}\\
                \boxed{\mathrm{c}_2} & \boxed{\mathrm{c}_3}\\
                \mathrm{} & \mathrm{}\\
                \mathrm{} & \mathrm{}
            \end{pmatrix}
        }_{\mathrm{c}_o}
        \\&=
        \begin{pmatrix}
            \pmb{\text{\Large }} & \pmb{\text{\Large }} & \pmb{\text{\Large }}\\
            \boxed{\mathrm{c}_i} & \boxed{\mathrm{v}_t} & \boxed{\mathrm{c}_o}\\
            \mathrm{} & \mathrm{} & \mathrm{}\\
            \mathrm{} & \mathrm{} & \mathrm{}
        \end{pmatrix}
        =
        \begin{pmatrix}\pmb{\text{\Large }}\\\boxed{\mathrm{syllable}}\\\mathrm{}\\\mathrm{}\end{pmatrix}_{\text{\footnotesize Answer: won}}
    \end{align}\\

    \begin{align}
        \textbf{\Huge 괜} &= 
        \underbrace{
            \begin{pmatrix}\pmb{\text{\Large }}\\\boxed{\mathrm{c}_1}\\ \mathrm{}\\ \mathrm{} \end{pmatrix}
        }_{\mathrm{c}_i}
        +
        \underbrace{
            \begin{pmatrix}             
                \pmb{\text{\Large }}&\pmb{\text{\Large }}\\
                \boxed{\mathrm{v}_1}&\boxed{\mathrm{v}_2}\\
                \mathrm{}&\mathrm{}\\
                \mathrm{}&\mathrm{}
            \end{pmatrix}
        }_{\mathrm{v}_t}
        +
        \underbrace{
            \begin{pmatrix}
                \pmb{\text{\Large }} & \pmb{\text{\Large }}\\
                \boxed{\mathrm{c}_2} & \boxed{\mathrm{c}_3}\\
                \mathrm{} & \mathrm{}\\
                \mathrm{} & \mathrm{}
            \end{pmatrix}
        }_{\mathrm{c}_o}
        \\&=
        \begin{pmatrix}
            \pmb{\text{\Large }} & \pmb{\text{\Large }} & \pmb{\text{\Large }}\\
            \boxed{\mathrm{c}_i} & \boxed{\mathrm{v}_t} & \boxed{\mathrm{c}_o}\\
            \mathrm{} & \mathrm{} & \mathrm{}\\
            \mathrm{} & \mathrm{} & \mathrm{}
        \end{pmatrix}
        =
        \begin{pmatrix}\pmb{\text{\Large }}\\\boxed{\mathrm{syllable}}\\\mathrm{}\\\mathrm{}\end{pmatrix}_{\text{\footnotesize Answer: gwaen}}
    \end{align}\\
    
    \begin{align}
        \textbf{\Huge 봤} &= 
        \underbrace{
            \begin{pmatrix}\pmb{\text{\Large }}\\\boxed{\mathrm{c}_1}\\ \mathrm{}\\ \mathrm{} \end{pmatrix}
        }_{\mathrm{c}_i}
        +
        \underbrace{
            \begin{pmatrix}             
                \pmb{\text{\Large }}&\pmb{\text{\Large }}\\
                \boxed{\mathrm{v}_1}&\boxed{\mathrm{v}_2}\\
                \mathrm{}&\mathrm{}\\
                \mathrm{}&\mathrm{}
            \end{pmatrix}
        }_{\mathrm{v}_t}
        +
        \underbrace{
            \begin{pmatrix}
                \pmb{\text{\Large }} & \pmb{\text{\Large }}\\
                \boxed{\mathrm{c}_2} & \boxed{\mathrm{c}_3}\\
                \mathrm{} & \mathrm{}\\
                \mathrm{} & \mathrm{}
            \end{pmatrix}
        }_{\mathrm{c}_o}
        \\&=
        \begin{pmatrix}
            \pmb{\text{\Large }} & \pmb{\text{\Large }} & \pmb{\text{\Large }}\\
            \boxed{\mathrm{c}_i} & \boxed{\mathrm{v}_t} & \boxed{\mathrm{c}_o}\\
            \mathrm{} & \mathrm{} & \mathrm{}\\
            \mathrm{} & \mathrm{} & \mathrm{}
        \end{pmatrix}
        =
        \begin{pmatrix}\pmb{\text{\Large }}\\\boxed{\mathrm{syllable}}\\\mathrm{}\\\mathrm{}\end{pmatrix}_{\text{\footnotesize Answer: bwadd}}
    \end{align}\\

    \begin{align}
        \textbf{\Huge 읽} &= 
        \underbrace{
            \begin{pmatrix}\pmb{\text{\Large }}\\\boxed{\mathrm{c}_1}\\ \mathrm{}\\ \mathrm{} \end{pmatrix}
        }_{\mathrm{c}_i}
        +
        \underbrace{
            \begin{pmatrix}             
                \pmb{\text{\Large }}&\pmb{\text{\Large }}\\
                \boxed{\mathrm{v}_1}&\boxed{\mathrm{v}_2}\\
                \mathrm{}&\mathrm{}\\
                \mathrm{}&\mathrm{}
            \end{pmatrix}
        }_{\mathrm{v}_t}
        +
        \underbrace{
            \begin{pmatrix}
                \pmb{\text{\Large }} & \pmb{\text{\Large }}\\
                \boxed{\mathrm{c}_2} & \boxed{\mathrm{c}_3}\\
                \mathrm{} & \mathrm{}\\
                \mathrm{} & \mathrm{}
            \end{pmatrix}
        }_{\mathrm{c}_o}
        \\&=
        \begin{pmatrix}
            \pmb{\text{\Large }} & \pmb{\text{\Large }} & \pmb{\text{\Large }}\\
            \boxed{\mathrm{c}_i} & \boxed{\mathrm{v}_t} & \boxed{\mathrm{c}_o}\\
            \mathrm{} & \mathrm{} & \mathrm{}\\
            \mathrm{} & \mathrm{} & \mathrm{}
        \end{pmatrix}
        =
        \begin{pmatrix}\pmb{\text{\Large }}\\\boxed{\mathrm{syllable}}\\\mathrm{}\\\mathrm{}\end{pmatrix}_{\text{\footnotesize Answer: ilg}}
    \end{align}


    \subsection{With special cases}
    \begin{align}
        \underbrace{\text{\Huge 몰}}_{\text{mol}} \underbrace{\textbf{\Huge 라}}_{\text{??}} &\overset{\text{special case}}{\underset{\text{\cref{eq:rieut}}}{=}} 
        \underbrace{
            \begin{pmatrix}\pmb{\text{\Large }}\\\boxed{\mathrm{c}_1}\\ \mathrm{}\\ \mathrm{} \end{pmatrix}
        }_{\mathrm{c}_i}
        +
        \underbrace{
            \begin{pmatrix}             
                \pmb{\text{\Large }}&\pmb{\text{\Large }}\\
                \boxed{\mathrm{v}_1}&\boxed{\mathrm{v}_2}\\
                \mathrm{}&\mathrm{}\\
                \mathrm{}&\mathrm{}
            \end{pmatrix}
        }_{\mathrm{v}_t}
        +
        \underbrace{
            \begin{pmatrix}
                \pmb{\text{\Large }} & \pmb{\text{\Large }}\\
                \boxed{\mathrm{c}_2} & \boxed{\mathrm{c}_3}\\
                \mathrm{} & \mathrm{}\\
                \mathrm{} & \mathrm{}
            \end{pmatrix}
        }_{\mathrm{c}_o}
        \\&=
        \begin{pmatrix}
            \pmb{\text{\Large }} & \pmb{\text{\Large }} & \pmb{\text{\Large }}\\
            \boxed{\mathrm{c}_i} & \boxed{\mathrm{v}_t} & \boxed{\mathrm{c}_o}\\
            \mathrm{} & \mathrm{} & \mathrm{}\\
            \mathrm{} & \mathrm{} & \mathrm{}
        \end{pmatrix}
        =
        \begin{pmatrix}\pmb{\text{\Large }}\\\boxed{\mathrm{syllable}}\\\mathrm{}\\\mathrm{}\end{pmatrix}_{\text{\footnotesize Answer: la}}
    \end{align}\\

    \begin{align}
        \underbrace{\textbf{\Huge 맞}}_{\text{??}} \underbrace{\text{\Huge 아}}_{\text{a}} &\overset{\text{special case}}{\underset{\text{\cref{eq:jieut}}}{=}} 
        \underbrace{
            \begin{pmatrix}\pmb{\text{\Large }}\\\boxed{\mathrm{c}_1}\\ \mathrm{}\\ \mathrm{} \end{pmatrix}
        }_{\mathrm{c}_i}
        +
        \underbrace{
            \begin{pmatrix}             
                \pmb{\text{\Large }}&\pmb{\text{\Large }}\\
                \boxed{\mathrm{v}_1}&\boxed{\mathrm{v}_2}\\
                \mathrm{}&\mathrm{}\\
                \mathrm{}&\mathrm{}
            \end{pmatrix}
        }_{\mathrm{v}_t}
        +
        \underbrace{
            \begin{pmatrix}
                \pmb{\text{\Large }} & \pmb{\text{\Large }}\\
                \boxed{\mathrm{c}_2} & \boxed{\mathrm{c}_3}\\
                \mathrm{} & \mathrm{}\\
                \mathrm{} & \mathrm{}
            \end{pmatrix}
        }_{\mathrm{c}_o}
        \\&=
        \begin{pmatrix}
            \pmb{\text{\Large }} & \pmb{\text{\Large }} & \pmb{\text{\Large }}\\
            \boxed{\mathrm{c}_i} & \boxed{\mathrm{v}_t} & \boxed{\mathrm{c}_o}\\
            \mathrm{} & \mathrm{} & \mathrm{}\\
            \mathrm{} & \mathrm{} & \mathrm{}
        \end{pmatrix}
        =
        \begin{pmatrix}\pmb{\text{\Large }}\\\boxed{\mathrm{syllable}}\\\mathrm{}\\\mathrm{}\end{pmatrix}_{\text{\footnotesize Answer: maj}}
    \end{align}\\

    \begin{align}
        \underbrace{\textbf{\Huge 옷}}_{\text{??}} \underbrace{\text{\Huge 이}}_{\text{i}} &\overset{\text{special case}}{\underset{\text{\cref{eq:shieut}}}{=}} 
        \underbrace{
            \begin{pmatrix}\pmb{\text{\Large }}\\\boxed{\mathrm{c}_1}\\ \mathrm{}\\ \mathrm{} \end{pmatrix}
        }_{\mathrm{c}_i}
        +
        \underbrace{
            \begin{pmatrix}             
                \pmb{\text{\Large }}&\pmb{\text{\Large }}\\
                \boxed{\mathrm{v}_1}&\boxed{\mathrm{v}_2}\\
                \mathrm{}&\mathrm{}\\
                \mathrm{}&\mathrm{}
            \end{pmatrix}
        }_{\mathrm{v}_t}
        +
        \underbrace{
            \begin{pmatrix}
                \pmb{\text{\Large }} & \pmb{\text{\Large }}\\
                \boxed{\mathrm{c}_2} & \boxed{\mathrm{c}_3}\\
                \mathrm{} & \mathrm{}\\
                \mathrm{} & \mathrm{}
            \end{pmatrix}
        }_{\mathrm{c}_o}
        \\&=
        \begin{pmatrix}
            \pmb{\text{\Large }} & \pmb{\text{\Large }} & \pmb{\text{\Large }}\\
            \boxed{\mathrm{c}_i} & \boxed{\mathrm{v}_t} & \boxed{\mathrm{c}_o}\\
            \mathrm{} & \mathrm{} & \mathrm{}\\
            \mathrm{} & \mathrm{} & \mathrm{}
        \end{pmatrix}
        =
        \begin{pmatrix}\pmb{\text{\Large }}\\\boxed{\mathrm{syllable}}\\\mathrm{}\\\mathrm{}\end{pmatrix}_{\text{\footnotesize Answer: os}}
    \end{align}
    \subsection{Whole words}
    \begin{align}
        \underbrace{\textbf{\Huge 우유}}_{\text{milk}} =
        \begin{pmatrix}
            \pmb{\text{\Large }} & \pmb{\text{\Large }}\\
            \boxed{\mathrm{syllable}_1} & \boxed{\mathrm{syllable}_2}\\
            \mathrm{} & \mathrm{}\\
            \mathrm{} & \mathrm{}
        \end{pmatrix}_{\text{\footnotesize Answer: oo-you}}
    \end{align}\\

     \begin{align}
        \underbrace{\textbf{\Huge 사랑}}_{\text{love}} =
        \begin{pmatrix}
            \pmb{\text{\Large }} & \pmb{\text{\Large }}\\
            \boxed{\mathrm{syllable}_1} & \boxed{\mathrm{syllable}_2}\\
            \mathrm{} & \mathrm{}\\
            \mathrm{} & \mathrm{}
        \end{pmatrix}_{\text{\footnotesize Answer: sa-rang}}
    \end{align}\\

     \begin{align}
        \underbrace{\textbf{\Huge 친구}}_{\text{friend}} =
        \begin{pmatrix}
            \pmb{\text{\Large }} & \pmb{\text{\Large }}\\
            \boxed{\mathrm{syllable}_1} & \boxed{\mathrm{syllable}_2}\\
            \mathrm{} & \mathrm{}\\
            \mathrm{} & \mathrm{}
        \end{pmatrix}_{\text{\footnotesize Answer: chin-goo}}
    \end{align}\\

     \begin{align}
        \underbrace{\textbf{\Huge 김치}}_{\text{fermented cabbage}} =
        \begin{pmatrix}
            \pmb{\text{\Large }} & \pmb{\text{\Large }}\\
            \boxed{\mathrm{syllable}_1} & \boxed{\mathrm{syllable}_2}\\
            \mathrm{} & \mathrm{}\\
            \mathrm{} & \mathrm{}
        \end{pmatrix}_{\text{\footnotesize Answer: gim-chi}}
    \end{align}\\

    \begin{align}
        \underbrace{\textbf{\Huge 한국}}_{\text{Korea}} =
        \begin{pmatrix}
            \pmb{\text{\Large }} & \pmb{\text{\Large }}\\
            \boxed{\mathrm{syllable}_1} & \boxed{\mathrm{syllable}_2}\\
            \mathrm{} & \mathrm{}\\
            \mathrm{} & \mathrm{}
        \end{pmatrix}_{\text{\footnotesize Answer: han-goog}}
    \end{align}\\

     \begin{align}
        \underbrace{\textbf{\Huge 신발}}_{\text{shoe}} =
        \begin{pmatrix}
            \pmb{\text{\Large }} & \pmb{\text{\Large }}\\
            \boxed{\mathrm{syllable}_1} & \boxed{\mathrm{syllable}_2}\\
            \mathrm{} & \mathrm{}\\
            \mathrm{} & \mathrm{}
        \end{pmatrix}_{\text{\footnotesize Answer: shin-bal}}
    \end{align}
    

\nocite{*}
\printbibliography
\end{document}